\documentclass[12pt,a4paper]{article}

\usepackage[a4paper,margin=1cm]{geometry}
\usepackage{minted}
\usepackage{enumitem}
\usepackage{setspace}

% \setstretch{1.15}
\setminted{
    fontsize=\small,
    breaklines,
    frame=single,
    framesep=2mm,
    baselinestretch=1
}

\pagenumbering{gobble}

\begin{document}

\begin{center}
  \textbf{KUIS 2}\\
  \textbf{ALGORITMA PEMROGRAMAN 2}
\end{center}

\begin{enumerate}[label=\arabic*.]
  \item Apa tampilan dari coding berikut ini (20 point):

        \begin{minted}{java}
public class Soal1 {
    public static void main(String args[]) {
        Hewan hewan1 = new kucing();
        System.out.println("Jumlah Hewan1 =" + hewan1.jumlah);
        hewan1.suara();

        harimau tiger = new harimau();
        System.out.println("Jumlah tiger =" + tiger.jumlah);
        tiger.suara(tiger);

        Hewan hewan2 = new harimau();
        System.out.println(" Jumlah Hewan2 =" + hewan2.jumlah);
        hewan2.suara();
    }
}

class Hewan {
    int jumlah = 0;

    public void suara() {
        System.out.println("....");
    }
}

class kucing extends Hewan {
    int jumlah = 10;

    public void info() {
        System.out.println("Puzz puzz");
    }
}

class harimau extends Hewan {
    int jumlah = 7;

    public void suara() {
        System.out.println("auuuu auung");
    }

    public void suara(harimau T) {
        System.out.println("guuk gukK");
    }
}
\end{minted}

  \item Lengkap coding berikut (20 point):

        \begin{minted}{java}
public class Soal2 {
    public static void main(String[] args) {
        Persegi A = new Persegi("Merah", 8);
        A.info();
    }
}

class Bentuk {
    String Warna;

    public Bentuk(String x) {
        System.out.println("Konstruktor Bentuk dipanggil");
        ................................
    }
}

class Persegi extends Bentuk {
    private int sisi;

    public Persegi(String A, int B) {
        ................
        System.out.println("Konstruktor Persegi dipanggil");
        ................
    }

    public void info() {
        System.out.println("Warna persegi adalah : " + Warna);
        System.out.println("Ukuran adalah : " + sisi);
    }
}
\end{minted}

        Agar menampilkan output berikut:

        \begin{minted}{text}
Konstruktor Bentuk dipanggil
Konstruktor Persegi dipanggil
Warna persegi adalah : Merah
Ukuran adalah : 8
\end{minted}

  \item Buatlah coding pada problem berikut:

        Seorang penjual alat tulis menjual 10 bolpoint, 10 pensil dan 10 penghapus.
        1 biji bolpoint harganya Rp. 2000, 1 biji pensil harganya Rp. 1.000 dan
        1 penghapus harganya Rp. 500. Gunakanlah objek untuk menyelesaikan soal
        di bawah ini!

        \begin{enumerate}[label=\alph*.]
          \item Buatlah method untuk memasukkan (\textit{setter}) nama, stok, harga satuan,
                dan harga (stok $\times$ harga satuan) alat tulis tersebut! (20 point)
          \item Buatlah method untuk menampilkan (\textit{getter}) nama, stok, harga satuan,
                dan harga (stok $\times$ harga satuan) alat tulis tersebut! (20 point)
          \item Buatlah method Total Harga (\textit{setter getter}) untuk menampilkan uang yang
                diterima penjual jika semua alat tulis tersebut terjual semuanya! (20 point)
        \end{enumerate}

        Agar menampilkan output:

        \begin{minted}{text}
Nama Alat Tulis : Bolpoint
Stok : 10
Harga Satuan: 2000
Harga Bolpoint : 20000
Nama Alat Tulis : Pensil
Stok : 10
Harga Satuan: 1000
Harga Pensil : 10000
Nama Alat Tulis : Penghapus
Stok : 10
Harga Satuan: 500
Harga Penghapus : 5000
Total Harga: 35000
\end{minted}
\end{enumerate}

\newpage

\begin{center}
  \textbf{SOLUSI KUIS 2}\\
  \textbf{ALGORITMA PEMROGRAMAN 2}
\end{center}

\begin{enumerate}[label=\arabic*.]
  \item Tampilan output dari program:

        \begin{minted}{text}
Jumlah Hewan1 =0
....
Jumlah tiger =7
guuk gukK
 Jumlah Hewan2 =0
auuuu auung
\end{minted}

        Penjelasan singkat: variabel \texttt{jumlah} mengikuti tipe referensi objek,
        sedangkan method \texttt{suara()} mengikuti method yang dioverride oleh objek
        sebenarnya.

  \item Coding lengkap:

        \begin{minted}{java}
public class Soal2 {
    public static void main(String[] args) {
        Persegi A = new Persegi("Merah", 8);
        A.info();
    }
}

class Bentuk {
    String Warna;

    public Bentuk(String x) {
        System.out.println("Konstruktor Bentuk dipanggil");
        Warna = x;
    }
}

class Persegi extends Bentuk {
    private int sisi;

    public Persegi(String A, int B) {
        super(A);
        System.out.println("Konstruktor Persegi dipanggil");
        sisi = B;
    }

    public void info() {
        System.out.println("Warna persegi adalah : " + Warna);
        System.out.println("Ukuran adalah : " + sisi);
    }
}
\end{minted}

        Output:

        \begin{minted}{text}
Konstruktor Bentuk dipanggil
Konstruktor Persegi dipanggil
Warna persegi adalah : Merah
Ukuran adalah : 8
\end{minted}

  \item Coding penyelesaian:

        \begin{minted}{java}
class AlatTulis {
    private String nama;
    private int stok;
    private int hargaSatuan;
    private int harga;
    private static int totalHarga;

    public void setNama(String nama) {
        this.nama = nama;
    }

    public void setStok(int stok) {
        this.stok = stok;
    }

    public void setHargaSatuan(int hargaSatuan) {
        this.hargaSatuan = hargaSatuan;
    }

    public void setHarga() {
        harga = stok * hargaSatuan;
        totalHarga += harga;
    }

    public String getNama() {
        return nama;
    }

    public int getStok() {
        return stok;
    }

    public int getHargaSatuan() {
        return hargaSatuan;
    }

    public int getHarga() {
        return harga;
    }

    public static int getTotalHarga() {
        return totalHarga;
    }

    public void tampilkanData() {
        System.out.println("Nama Alat Tulis : " + getNama());
        System.out.println("Stok : " + getStok());
        System.out.println("Harga Satuan: " + getHargaSatuan());
        System.out.println("Harga " + getNama() + " : " + getHarga());
    }
}

public class Soal3 {
    public static void main(String[] args) {
        AlatTulis bolpoint = new AlatTulis();
        bolpoint.setNama("Bolpoint");
        bolpoint.setStok(10);
        bolpoint.setHargaSatuan(2000);
        bolpoint.setHarga();

        AlatTulis pensil = new AlatTulis();
        pensil.setNama("Pensil");
        pensil.setStok(10);
        pensil.setHargaSatuan(1000);
        pensil.setHarga();

        AlatTulis penghapus = new AlatTulis();
        penghapus.setNama("Penghapus");
        penghapus.setStok(10);
        penghapus.setHargaSatuan(500);
        penghapus.setHarga();

        bolpoint.tampilkanData();
        pensil.tampilkanData();
        penghapus.tampilkanData();

        System.out.println("Total Harga: " + AlatTulis.getTotalHarga());
    }
}
\end{minted}

        Output:

        \begin{minted}{text}
Nama Alat Tulis : Bolpoint
Stok : 10
Harga Satuan: 2000
Harga Bolpoint : 20000
Nama Alat Tulis : Pensil
Stok : 10
Harga Satuan: 1000
Harga Pensil : 10000
Nama Alat Tulis : Penghapus
Stok : 10
Harga Satuan: 500
Harga Penghapus : 5000
Total Harga: 35000
\end{minted}
\end{enumerate}

\end{document}
